\section{A precise route toward quasi-polynomial complexity}
\label{sec:roadmap}

The new arithmetic stage is polynomial on its supplied input class.
The bounded-description search is quasi-polynomial on a restricted
one-sided discovery problem. Neither result bounds the cost of the
remaining general instances. The objective now is to identify an interface
and a discovery theorem that can survive those instances without
materializing the exponential objects isolated above.

\subsection{Which compression target should be changed}

There are at least four different mathematical contracts for a tangle
summary:
\begin{enumerate}
\item \emph{Relative-boundary isotopy information.}
A rational slope determines a rational tangle up to isotopy fixing its
marked boundary. This is exact topological information, compactly encoded.
\item \emph{Full homotopy information.}
A Bar--Natan complex retains enough data to compute the homology of every
planar composition. Explicit representatives can be exponentially large
at the prescribed checkpoints of Section~\ref{sec:unknot-checkpoint-lower-bound}.
\item \emph{All Boolean continuation answers.}
The relation \(P\sim P'\) means that every allowed continuation produces
the same unknot decision. Section~\ref{sec:rational-continuation-rank}
shows exponentially many classes even for a fixed four-point boundary.
It also rules out a small linear factorization of that final Boolean
predicate on the constructed family.
\item \emph{A certificate for the actual input.}
A Montesinos arithmetic proof or a forbidden local disk can decide
one input without constructing all continuation answers. This is the
contract used by the delivered early exits.
\end{enumerate}

These contracts should be reflected in types and function boundaries.
A local nontriviality witness cannot be inserted as a substitute for a
full chain complex at an arbitrary internal node. A rational slope can
represent the topology of a rational piece, but this alone does not provide
a fast operation for gluing it to an arbitrary nonrational complex.
The required operation must either have its own theorem or retain the
original diagram for an exact fallback.

The rank-two identity
\[
 D_m=ps^{\mathsf T}-qr^{\mathsf T},
 \qquad H_m=\mathbf1\{|D_m|=1\}
\]
illustrates a productive distinction. A low-dimensional arithmetic value
can encode an exact answer after a nonlinear readout even when the
Boolean connection matrix has full rank. State-count and Boolean-rank
obstructions should therefore be used to reject an unsuitable
representation contract, not to reject all arithmetic approaches.

\subsection{What the proved restricted discovery theorem actually buys}

Proposition~\ref{prop:mon:bounded-discovery} removes the need for a
supplied occurrence map on a specified family. It is a genuine algorithm:
enumerate bounded sources, reject arithmetic or component failures,
compile a literal plane pattern, and propagate each possible root image.
When \(h=O(\log n)\), its running time is \(n^{O(\log n)}\), including
bit arithmetic and the polynomial work per candidate.

The quantifiers are the important part. Completeness holds for forbidden
tangles that already appear as literal admissible diagrams with the
specified source bound, orientation, and four distinct external edges.
The algorithm does not search arbitrary isotopies. It does not prove
that every nontrivial diagram contains a short witness of this kind.
It also never certifies an unknot merely because no forbidden template
was found.

The hypothesis is substantial. A continued fraction with a long sequence
of coefficients equal to two can have \(h=\Theta(n)\); storing its slope
still takes only \(O(n)\) bits, but the enumeration theorem's logarithmic
coefficient-count hypothesis need not hold. Conversely, a large twist
coefficient can contribute many displayed crossings while occupying one
coefficient slot. This explains why coefficient count is useful and why
it is not a universal measure of diagram difficulty.

The fixed catalogue already gives a concrete scalable rejection family
with arbitrary exteriors: every accepted literal occurrence of
\(R(1/3)+R(1/4)\), for example, proves nontriviality regardless of the
number or complexity of exterior crossings. A useful next theorem would
enlarge this family through certified local rewrites or a rational-region
extractor while preserving a proved total discovery bound.

\subsection{A hierarchy accounting criterion}

The geometric hierarchy route is complementary to the arithmetic tangle
route. It should be audited using the iteration count and operation cost
as separate parameters.

\begin{proposition}[Conditional accounting, not a new hierarchy theorem]
\label{prop:hierarchy-accounting}
Suppose a fully specified implementation has at most \(L(g+1)^L\)
elementary hierarchy events, as in the iteration estimate of
\cite[Proposition~9.1]{lackenby-2026}. Suppose additionally that
\[
 g\le n^a,\qquad L\le b(\log(n+2))^c
\]
for fixed constants \(a,b,c>0\), and that all encoding growth and
bit operations at each event are bounded by
\(\exp(O((\log(n+2))^{c+1}))\).
Then its total bit time has the same form
\[
 \exp(O((\log(n+2))^{c+1})).
\]
In particular, \(c=1\) yields \(n^{O(\log n)}\), whereas \(c=2\)
yields \(n^{O((\log n)^2)}\).
\end{proposition}

\begin{proof}
Taking logarithms of the event bound gives
\[
 \log L+L\log(g+1)
 =O((\log(n+2))^{c+1}).
\]
Multiplying the number of events by the assumed per-event bound adds
two quantities of that order in the exponent. This accounts for both
iteration count and operation cost. It proves no one of the hypotheses.
\end{proof}

This accounting is useful precisely because each hypothesis can fail
independently. A short hierarchy does not ensure short coordinate
descriptions. A combinatorially decreasing potential does not bound the
cost of finding the next compression. Replacing a surface by a compressed
description does not make intersection and update operations efficient
without an explicit algorithm. These obligations are central to turning
the hierarchy programme into a reviewed implementation.

The present tangle lower bounds do not refute such a geometric approach.
They concern specific algebraic summaries. Their practical lesson is
that geometric events should not be forced through an explicit full
complex merely to preserve an interface whose information is richer
than the decision requires.

\section{Further research questions and concrete milestones}
\label{sec:questions}

The following questions are ordered roughly from immediately actionable
implementation work to broader mathematical objectives. Each specifies
what would count as progress. They are research proposals, not assumptions
used to justify the existing code.

\begin{researchquestion}[Certified extraction of visible rational regions]
Can a validated PD be reduced to maximal visible rational regions in
\(O(n\log n)\) or another explicit polynomial time, while returning
relative-boundary isotopy certificates and exact slopes?
\end{researchquestion}
Start with bigon chains and the inverse operations of the rational compiler,
retaining cyclic port order, crossing signs, and a replayable contraction
tree. A useful first result would recognize the supplied diagonal family
after source metadata is removed and after arbitrary relabeling and
crossing reordering. A stronger result would admit carefully specified
flypes or local isotopies. The theorem should define which diagrams
the extractor is complete on, and charge each graph update rather than
assuming repeated scans are free.

\begin{researchquestion}[Choosing the placement of local search]
Can local obstruction search run after all cheaper certificates and before
the expensive fallback with a small, provable overhead bound?
\end{researchquestion}
The current negative controls already demonstrate why this matters.
Measure construction, validation, compilation, discovery, and fallback
separately, using held-out hosts as well as planted patterns. Share target
dart arrays and precompiled templates. A useful release gate is an
end-to-end improvement on difficult plain PDs while bounding overhead on
trefoils, descending diagrams, and homogeneous examples. A benchmark
selected only from successful local witnesses would not answer the question.

\begin{researchquestion}[Practical bounded-description enumeration]
Can the \(n^{O(h)}\) enumeration be pruned substantially without losing
completeness for its stated literal class?
\end{researchquestion}
Possible exact pruning keys include expansion size, endpoint pairing,
crossing-sign counts, rooted ribbon neighborhoods, and arithmetic residues.
Canonicalization must respect the difference between two isotopic rational
sources and two different literal diagrams: equality of slopes is enough
for isotopy, but not necessarily enough for an unmodified dart match.
An implementation should return the same existence answer as the exhaustive
small-\(h\) reference enumerator before any asymptotic speed claim is made.

\begin{researchquestion}[A larger certified local-tangle language]
Which arborescent or graph-manifold tangle classes admit compact peripheral
data and complete tests for whether an unknotting complement can exist?
\end{researchquestion}
The Montesinos criterion uses more than homology order: its congruence
condition excludes \(R(1/3)+R(1/4)\) even though the numerator and denominator
determinants are coprime. An extension should retain the exact boundary
slopes, exceptional-fiber information, and marking needed for gluing.
It should classify rational and product degeneracies explicitly. Merely
recording the first homology group of each piece loses the peripheral
information that makes the current obstruction effective.

\begin{researchquestion}[Typed arithmetic composition with an arbitrary exterior]
What exact interface permits a rational slope or a compact Seifert piece
to interact with a nonrational remainder without expanding its full complex?
\end{researchquestion}
A first milestone could support one rational piece glued to a small explicit
remainder with a polynomial cost in its bit description and the remainder's
actual size. The correctness theorem must specify which topological or
homological answer the interface preserves. There should be explicit
conversion and fallback rules. A compact slope object is not automatically
an object of the pre-existing chain-complex category.

\begin{researchquestion}[Continuation-dependent processing of the Pell family]
Can a general PD strategy detect and exploit the relationship between
prefix and actual suffix before it reaches the exponential checkpoint?
\end{researchquestion}
On the diagonal family, the decisive relation is the unimodular pairing
\(p_ks_k-q_kr_k=-1\). An algorithm that deliberately constructs the
full prefix complex pays for information that the suffix immediately
collapses. A useful experiment would compare symmetric gluing,
look-ahead, and arithmetic extraction on source-erased diagrams with
controlled perturbations. A useful theorem would bound the size of all
intermediate representations for that algorithm, rather than only
bounding the final unknot rank.

\begin{researchquestion}[Compact homological data beyond multiplicities]
Can differentials and extension data for structured tangle complexes be
stored in an arithmetic circuit, recurrence, or other succinct form on
which the needed closure operations remain efficient?
\end{researchquestion}
The lower bound permits such representations because it counts expanded
summands. It does not prove they exist with efficient operations.
A scalar multiplicity alone omits differential data; capping a multiplicity
can be unsound if later maps interact with its copies. The target should
include composition, cancellation, restriction to a justified homological
window, and exact rank decisions on the permitted output, with explicit
bit-growth bounds. Symbolic expression size must be analyzed alongside
the number of graph nodes in the representation.

\begin{researchquestion}[Useful arithmetic values before a nonlinear readout]
Which intermediate arithmetic invariants have a low-dimensional gluing
law but a sufficiently strong exact terminal predicate?
\end{researchquestion}
The signed determinant pairing is a complete example on rational
continuations. It converts a rank-two arithmetic matrix into a full-rank
Boolean matrix by testing membership in \(\{1,-1\}\).
For more general pieces, one might study peripheral matrices, presentations
with controlled relations, or exact equations for distinguished fillings.
An arithmetic invariant that is not complete should remain a one-sided
filter. In particular, a compact Jones or Euler calculation should not be
promoted to a general unknot detector without an additional theorem.

\begin{researchquestion}[Generalized representation lower bounds]
Can the continuation construction distinguish other proposed summary
models, including bounded-degree nonlinear readouts, compressed modules,
or restricted arithmetic circuits?
\end{researchquestion}
Any proposed lower bound must specify the model and its input encoding.
The current identity matrix already implies that a representative-subfamily
scheme preserving the existence of an accepting prefix for every displayed
suffix must retain all \(2^m\) prefixes: suffix \(Q_w\) accepts only \(P_w\).
This does not prevent a symbolic set representation. A stronger useful
result would identify exactly which operations force expansion and which
ones preserve compactness, rather than claiming a general complexity
separation from a state count.

\begin{researchquestion}[Sharper, simpler independent verification]
Can the separate Montesinos verifier be reduced from the conservative
cubic bound to \(O(B^2)\), or better with fast arithmetic, without merely
reusing the production implementation?
\end{researchquestion}
The nested replay exists to check tracked prefix statistics, not to decide
the topology. One option is to remove nonessential statistics from the
mathematical certificate and verify them separately. Another is to maintain
independent prefix products or use a product tree. The important outcome is
an independent small checker with a transparent source-to-fraction and
fraction-to-verdict proof, not a second function that trusts the producer.

\begin{researchquestion}[Short positive certificates and formal proof kernels]
Can a determinant-one arithmetic decision be turned into a compact,
replayable topological certificate and formalized in the ProveIt setting?
\end{researchquestion}
For compressed sources, a fully expanded Reidemeister movie may be longer
than the source by an exponential factor. A suitable certificate should
instead record rational-tangle identities, Euclidean arithmetic, and
boundary-preserving gluing steps in compressed form. Its checker must
verify all signs and primitive pairs. For local negative certificates,
the disk lemma can be formalized as a statement about finite dart
permutations and a plane embedding, with the imported topological theorem
kept explicit. Formalizing the disjoint-cone and triangle-representation
lemmas would give smaller initial proof-assistant milestones.

\begin{researchquestion}[Extending the local certificate geometry]
Can the implementation safely admit two ports on a single exterior edge,
disconnected crossing shadows with simple arcs, or additional boundary
marking conventions?
\end{researchquestion}
The initial four-distinct-edge contract avoids these cases and is sound.
An extension needs a new explicit cut-point representation and a proof
that the accepted common face still bounds the intended disk with no
hidden components. Disconnected shadows require more than simply
reusing root propagation from one component. Separate topology lemmas
and adversarial tests should precede any expansion of the accepted class.

\begin{researchquestion}[A useful structural completeness theorem]
Can every diagram in a broad, explicitly specified class be transformed
or decomposed so that arithmetic classification or a short forbidden
local certificate becomes available within quasi-polynomial cost?
\end{researchquestion}
Candidate classes should be defined by their diagrams and allowed moves,
not retrospectively by whether the implementation succeeds.
One could begin with a controlled arborescent grammar or diagrams with
a supplied decomposition tree and then study discovery.
Both the transformation count and the size of each intermediate diagram
or compressed source must be bounded. The current results provide
neither a universal \(O(\log n)\) coefficient bound nor a theorem that
local forbidden tangles cover all nontrivial knots.

\begin{researchquestion}[Hierarchy operation costs and the logarithmic exponent]
Can the geometric hierarchy programme supply executable bounds for
compression discovery, parallelity-bundle updates, boundary-pattern
simplification, and terminal checks, with controlled coordinate growth?
\end{researchquestion}
The accounting in Proposition~\ref{prop:hierarchy-accounting} identifies
the necessary parameters. An implementation audit should attach an
explicit representation and cost to each operation from the mathematical
algorithm. It should separately establish whether the needed hierarchy
length is \(O(\log n)\) or \(O((\log n)^2)\), since those produce
different quasi-polynomial exponents under polynomial pattern complexity.
A declining hierarchy potential alone is insufficient.

\begin{researchquestion}[Adversarial families and fair backend comparisons]
Which examples remain hard after provenance is removed, cheap certificates
are applied, and the best supported existing backend is selected?
\end{researchquestion}
The current diagonal family isolates a real expanded-complex problem;
the Alexander-one family isolates the failure of a common invariant;
the planted local hosts test certificate discovery. A broader study should
separate these purposes and include source-erased rational diagrams,
hidden local patterns, examples with no matches, and alternate compositions.
Keep resource exits censored, use fresh caches, report process startup
when relevant, and compare end-to-end time as well as inner kernels.
This would guide algorithm design without confusing favourable source
metadata with a general discovery improvement.

\subsection{Recommended sequence of follow-up contributions}

The next contribution should implement and certify visible rational-region
extraction, then benchmark the present arithmetic decision on ordinary
source-erased PDs. That would remove the largest practical qualification
from the supplied-source speedups on an explicit class.

A second contribution should move or schedule local search using the
measured pipeline costs, add an exact small-\(h\) reference enumerator,
and develop completeness-preserving pruning. Its claim should remain a
local-witness discovery bound until a wider structural theorem is proved.

A third contribution should implement one typed arithmetic/exterior
composition operation with a full correctness proof and a proved bound on
all intermediate representations. The Pell checkpoint family gives a sharp
test: success should avoid the expanded \(p_k+q_k\) cost, not simply hide
it in a later conversion step.

The broad quasi-polynomial objective then requires a coverage or hierarchy
theorem together with executable operation costs. The delivered arithmetic
and local certificate stages are exact components available for that
programme. The lower bounds identify a representation choice that can
defeat the programme even on unknots, while the arithmetic construction
demonstrates a concrete way around it in a mathematically complete case.
